\begin{corollary}[Conditional excess-mass tails]
With the notation of (PR1), for every $0<s<1$,
\[
 \mathbb E\left[e^{s\omega(P-P_{\min})}\mid\theta,A_\theta c=b\right]
 \le(1-s)^{-d_\theta}.
\]
Consequently, for every $t>0$, the conditional probability that
\[
 P-P_{\min}>\omega^{-1}\{d_\theta+2\sqrt{d_\theta t}+2t\}
\]
is at most $e^{-t}$.
\end{corollary}
\begin{proof}
Write $W(u)=V(u)-V(0)$. It is convex, nonnegative, and vanishes at
zero. For $r=(1-s)^{-1}\ge1$, convexity gives $W(ru)\ge rW(u)$.
A change of variables therefore shows
\[
 \int e^{-\omega W(u)/r}\,du
 =r^{d_\theta}\int e^{-\omega W(ru)/r}\,du
 \le r^{d_\theta}\int e^{-\omega W(u)}\,du.
\]
This is the stated moment bound. The standard Chernoff calculation
for $(1-s)^{-d_\theta}$ gives the displayed, slightly relaxed,
Gamma-tail threshold. If $d_\theta=0$, the excess is identically zero.
\end{proof}
